% Paste this environment into the IEEEtran source.
% The preamble must already \input the sibling iscafig.tex.
\begin{figure*}[t]
\centering
\begin{tikzpicture}[iscafig]
  \begin{scope}[local bounding box=map]
    \node[lbl, anchor=west] at (0,0.42) {\textbf{Caller frame}};
    \rec{0}{+0x00}{fldimm}{6.5cm}{previous frame}
    \rec{-0.52}{+0x04}{fldimm}{6.5cm}{raw code}
    \rec{-1.04}{+0x08}{fldimm}{6.5cm}{return PC}
    \rec{-1.56}{+0x0C}{fldimm}{6.5cm}{opcode PC}
    \rec{-2.08}{+0x10}{fldimm}{6.5cm}{flags, $n_{\mathrm{loc}}$, exception depth, $sp$}
    \rec{-2.60}{+0x14}{fldimm}{6.5cm}{continuation and return kind}
    \rec{-3.12}{+0x18}{fldmem}{6.5cm}{globals}
    \rec{-3.64}{+0x1C}{fldmem}{6.5cm}{\texttt{loc}}
    \rec{-4.16}{+0x20}{fldmem}{6.5cm}{function}
    \rec{-4.68}{+0x24}{fldimm}{6.5cm}{aux, aux2, aux3, through +0x2C}
    \rec{-5.20}{+0x30}{fldimm}{6.5cm}{qstr base; const base at +0x34}
    \rec{-5.72}{+0x38}{fldmem}{6.5cm}{generator}
    \rec{-6.24}{+0x3C}{fldimm}{6.5cm}{spill pointer}
    \rec{-6.76}{+0x40}{fldhot}{6.5cm}{spilled count and $sb$}
    \rec{-7.40}{+0x44}{fldcore}{6.5cm}{locals\,[0\,:\,$n_{\mathrm{loc}}$)}
    \rec{-8.04}{}{fldacc}{6.5cm}{exception slots, three words each}
    \rec{-8.68}{ftop}{fldhot}{6.5cm}{resident ring, logical order, $sp$ words}
    \node[lbl, anchor=west] at (0,-9.25)
      {callee frame base $=$ end of this snapshot};
  \end{scope}

  \begin{scope}[xshift=9.35cm]
    \node[lbl, anchor=west] at (0,0.42) {\textbf{On \texttt{CALL}}};
    \node[imm, anchor=north west, align=left, text width=7.3cm] (s1) at (0,-0.15)
      {\textbf{1.}\ Flush the local cache into the caller frame.};
    \node[imm, anchor=north west, align=left, text width=7.3cm, below=2.2mm of s1] (s2)
      {\textbf{2.}\ Write the 17-word header.
       The row at \texttt{+0x24} is three temporaries.
       The next row is the qstr base and the const base.};
    \node[imm, anchor=north west, align=left, text width=7.3cm, below=2.2mm of s2] (s3)
      {\textbf{3.}\ Store the resident ring at \texttt{ftop}, logical slot 0 first.
       An \texttt{RF\_LINK} in that ring still names the spill chunks.};
    \node[hot, anchor=north west, align=left, text width=7.3cm, below=2.2mm of s3] (s4)
      {\textbf{4.}\ The callee frame begins on the next word.
       Its ring is cleared: $sp=0$, $sb=0$.
       Locals are filled from the call, not left in the ring.};
    \node[core, anchor=north west, align=left, text width=7.3cm, below=2.2mm of s4] (s5)
      {\textbf{5.}\ Return copies that snapshot back into the ring and restores $sp$ and $sb$.};

    \node[mem, anchor=north west, align=left, text width=7.3cm] (bounds) at (0,-7.15)
      {\textbf{Frame reserve, 32\,KB}\\[1pt]
       \scriptsize \texttt{0x2000\_0000} frames grow upward\\
       \scriptsize \texttt{0x2000\_7E00} chunks grow downward\\
       \scriptsize a frame that would meet the chunks raises recursion};
  \end{scope}
\end{tikzpicture}
\caption{Caller frame written by a call.
The resident ring is snapshotted at the end of the caller, after the locals and the exception slots, and the callee is allocated immediately past that snapshot.
Spill chunks stay where they are.
The header field $sb$ is the physical register that holds logical slot 0.
Green words are pointers.
The rose words are the ring metadata the return path reloads.}
\label{fig:frame}
\end{figure*}
